% Prose rendering of PrimeTensor/Analysis/Chain.lean.
% Fragment: no \documentclass, no preamble, no document environment.

\section{The multiplicative chain rule}
\label{Analysis:Chain:sec}

\leanfile{PrimeTensor/Analysis/Chain.lean}

\subsection*{What this file establishes}

The chain rule, under hypotheses the module header is unusually explicit
about.  Differentiability of the two factors and scale control of their tangent
maps are, in its words, \emph{almost} sufficient; one further hypothesis is
exposed rather than hidden, namely that negligibility measured against the
actual inner response germ implies negligibility measured against the original
perturbation.  That is $\name{ChainAdmissibleAt}$, and
Design note~\ref{Analysis:Chain:note:why-hypothesis} is about why it has to be
assumed here when its classical counterpart is automatic.

Around it the file assembles: composition of tangent maps; the fact that a
differentiable map with a controlled tangent has a response that really does
approach the pivot, so that the inner response is itself a legitimate
perturbation; and an exact identity saying the response of a composite is the
outer response along the inner response.  The proof of the rule is then a
telescoping product of two ratios, and
Design note~\ref{Analysis:Chain:note:telescope} sets it beside the classical
two-term error split.

\subsection{Ratio algebra}

\begin{lemma}[\lean{MulReal.ratio\_comp}, \lean{ratio\_mul\_right} and
  \lean{mul\_ratio\_right}]\label{Analysis:Chain:lem:ratio}
$\name{ratio}(a,c) = \name{ratio}(a,b)\,\name{ratio}(b,c)$, and
$\name{ratio}(a,b) \cdot b = a = b \cdot \name{ratio}(a,b)$.
\end{lemma}

\begin{leancode}
theorem ratio_mul_right (a b : MulReal) :
    ratio a b * b = a := by
  unfold ratio
  calc
    (a * b⁻¹) * b = a * (b⁻¹ * b) := mul_assoc a b⁻¹ b
    _ = a * 1 := by rw [inv_mul b]
    _ = a := mul_one a
\end{leancode}

\begin{proof}
The first is the argument of \texttt{PrimeTensor/Scale/Laws.lean} one layer up:
reassociate to expose $b^{-1}b$, cancel, drop the unit.  The second is the same
cancellation with the intermediate magnitude restored, and the third is the
second after commuting.
\end{proof}

The cocycle identity and the reconstruction identity are the two facts the main
proof runs on.  The first telescopes a chain of comparisons; the second says
that a magnitude is recoverable from a comparison and its reference point,
which is what lets the inner response be treated as a perturbation in its own
right.

\begin{lemma}[\lean{MulReal.converges\_congr} and \lean{negligible\_congr}]
\label{Analysis:Chain:lem:congr}
Convergence and negligibility are unchanged by replacing a germ with one equal
to it term by term.
\end{lemma}

\begin{proof}
Keep the anchor and rewrite pointwise.
\end{proof}

These are bookkeeping, but they carry real weight below: several of the germs
in the main proof are equal on the nose rather than merely asymptotic, and
these two lemmas are how an exact rewriting is admitted into a statement about
levels.

\subsection{Composition of tangent maps}

\begin{definition}[\lean{MulTangentMap.comp}]\label{Analysis:Chain:def:comp}
$\name{comp}(D_g, D_f)(h) \defeq D_g(D_f(h))$, with the three morphism laws
proved by rewriting with the corresponding laws of the two factors, and
$\name{comp\_apply}$ recording the value by $\ctor{rfl}$.
\end{definition}

\begin{leancode}
def comp (Dg Df : MulTangentMap) : MulTangentMap where
  toFun := fun h => Dg (Df h)
  map_one := by
    rw [Df.map_one, Dg.map_one]
  map_mul := by
    intro a b
    rw [Df.map_mul, Dg.map_mul]
  map_inv := by
    intro a
    rw [Df.map_inv, Dg.map_inv]
\end{leancode}

\begin{lemma}[\lean{MulTangentMap.map\_ratio}]\label{Analysis:Chain:lem:map-ratio}
Every tangent map preserves ratios:
$D(\name{ratio}(a,b)) = \name{ratio}(D(a), D(b))$.
\end{lemma}

\begin{proof}
Unfold the ratio and apply preservation of products and of inverses.
\end{proof}

This is what lets an error, which is a ratio, be pushed through a tangent map
and still be recognised as an error between the images.

\subsection{Differentiability gives a genuine perturbation}

\begin{lemma}[\lean{negligible\_converges}]\label{Analysis:Chain:lem:neg-conv}
If $\name{error}$ is negligible relative to $\name{base}$ and $\name{base}$
approaches the pivot, then $\name{error}$ converges to the pivot.
\end{lemma}

\begin{proof}
Fix a level.  Take the anchor negligibility supplies at the smallest gain
$\ctor{one}$, and the anchor the convergence of the base supplies at the
requested level, and join them.  Beyond the join the base achieves that level,
so negligibility puts the error at $\name{advance}(\ell, \ctor{one})$, which is
$\ctor{succ}\,\ell$; one weakening by
\texttt{PrimeTensor/Analysis/Control.lean} gives $\ell$.
\end{proof}

\begin{theorem}[\lean{response\_converges}]\label{Analysis:Chain:thm:resp-conv}
If $f$ has derivative $D$ at $x$, $D$ is scale controlled, and $h$ approaches
the pivot, then $\name{response}(f,x,h)$ converges to the pivot.
\end{theorem}

\begin{proof}
The response is recovered from its error and the model by
Lemma~\ref{Analysis:Chain:lem:ratio}:
\[
  \name{response}_n
    = \name{ratio}\bigl(\name{response}_n,\, D(h_n)\bigr) \cdot D(h_n).
\]
The first factor converges to the pivot by
Lemma~\ref{Analysis:Chain:lem:neg-conv}, since differentiability makes it
negligible relative to $h$ and $h$ approaches the pivot.  The second converges
to the pivot by $\name{maps\_converges}$ of
\texttt{PrimeTensor/Analysis/Control.lean}, which is where control of $D$ is
used.  Their product converges to $1 \cdot 1 = 1$ by
\texttt{PrimeTensor/Analysis/Near/Laws.lean}, and
Lemma~\ref{Analysis:Chain:lem:congr} transports that along the pointwise
identity above.
\end{proof}

\begin{remark}[Differentiable implies continuous --- but not for free]
\label{Analysis:Chain:rem:diff-cont}
This is the statement that a differentiable map is continuous at the point, in
the only form the vocabulary offers: the response along any perturbation is
itself a perturbation.  In the additive theory it needs no extra hypothesis.
Here it needs $D$ to be scale controlled, and unavoidably so ---
\texttt{PrimeTensor/Analysis/Derivative.lean} placed no regularity on tangent
maps, so an unbounded model could carry a perturbation away from the pivot and
the response with it.  Control is what excludes that, and this theorem is the
first place the notion is spent rather than established.
\end{remark}

\subsection{The admissibility hypothesis}

\begin{definition}[\lean{ChainAdmissibleAt}]\label{Analysis:Chain:def:admissible}
$f$ is \emph{chain admissible at} $x$ when, for every $h$ approaching the
pivot, every germ negligible relative to $\name{response}(f,x,h)$ is negligible
relative to $h$.
\end{definition}

\begin{leancode}
def ChainAdmissibleAt
    (f : MulReal → MulReal) (x : MulReal) : Prop :=
  ∀ h : MulReal.Seq,
    ApproachesPivot h →
    ∀ error : MulReal.Seq,
      NegligibleRelative error (response f x h) →
      NegligibleRelative error h
\end{leancode}

\begin{designnote}[Why this is a hypothesis and not a lemma]
\label{Analysis:Chain:note:why-hypothesis}
The header identifies the classical counterpart: the comparison
$D f(h) = O(h)$ that transports an outer little-o error back to the original
variable.  Classically that comparison is not assumed --- for a differentiable
map with a bounded derivative it follows, since the increment is bounded by a
constant times the perturbation.

It cannot be lifted here in that form, for a reason
\texttt{PrimeTensor/Analysis/Derivative.lean} already recorded.  Comparisons in
this development are made with $\name{advance}$, which refines a level by a
\emph{positive} number of steps, there being no advance by zero.  So the
vocabulary can say ``strictly finer by a prescribed amount'' --- little-o ---
and cannot say ``no coarser than, up to a constant'' --- big-$O$.  The
statement the chain rule needs is a big-$O$ statement, and there is no way to
write it.

$\name{ChainAdmissibleAt}$ is what can be written instead.  Rather than
bounding the inner response against the perturbation, it asserts the
consequence that bounding would have: that the two germs are interchangeable as
\emph{bases} for negligibility, in the direction the proof consumes.  This
converts an inexpressible inequality into an expressible implication, at the
cost of making it an assumption.  Whether it follows from the other hypotheses
of Theorem~\ref{Analysis:Chain:thm:chain} is not settled in this file, and no
argument either way is given.
\end{designnote}

\begin{lemma}[\lean{chainAdmissibleAt\_identity}]
\label{Analysis:Chain:lem:adm-id}
The identity map is chain admissible at every point.
\end{lemma}

\begin{proof}
Its response along $h$ at stage $n$ is $\name{ratio}(x h_n, x)$.  That is
$h_n$, by the reconstruction identity of
\texttt{PrimeTensor/Analysis/Derivative.lean}.
So the two bases are the same germ and the implication is trivial: rewrite the
base and hand back the hypothesis.
\end{proof}

\subsection{Responses of a composite}

\begin{lemma}[\lean{response\_comp}]\label{Analysis:Chain:lem:resp-comp}
For every $n$,
\[
  \name{response}(g \circ f,\, x,\, h)_n
    = \name{response}\bigl(g,\, f(x),\, \name{response}(f,x,h)\bigr)_n .
\]
\end{lemma}

\begin{leancode}
theorem response_comp
    (f g : MulReal → MulReal)
    (x : MulReal)
    (h : MulReal.Seq) :
    ∀ n : Depth,
      response (fun y => g (f y)) x h n =
        response g (f x) (response f x h) n := by
  intro n
  unfold response
  rw [MulReal.mul_ratio_right
    (f (x * h n))
    (f x)]
\end{leancode}

\begin{proof}
The left side is $\name{ratio}\bigl(g(f(x h_n)),\, g(f(x))\bigr)$.  The right
side is $\name{ratio}\bigl(g(f(x) \cdot \name{inner}_n),\, g(f(x))\bigr)$ with
$\name{inner}_n = \name{ratio}(f(x h_n), f(x))$, and
Lemma~\ref{Analysis:Chain:lem:ratio} reconstructs
$f(x) \cdot \name{inner}_n = f(x h_n)$.  The two sides are then the same term.
\end{proof}

An identity, with no error and no hypothesis on $h$: the perturbation of the
outer function is exactly the response of the inner one.  This is what makes
the chain rule a matter of composing two first-order statements rather than of
estimating a composite directly.

\subsection{The chain rule}

\begin{theorem}[\lean{hasMulDerivativeAt\_comp}]\label{Analysis:Chain:thm:chain}
Suppose $f$ has derivative $D_f$ at $x$ with $D_f$ scale controlled, $f$ is
chain admissible at $x$, and $g$ has derivative $D_g$ at $f(x)$ with $D_g$
scale controlled.  Then $y \mapsto g(f(y))$ has derivative
$\name{comp}(D_g, D_f)$ at $x$.
\end{theorem}

\begin{leancode}
theorem hasMulDerivativeAt_comp
    {f g : MulReal → MulReal}
    {x : MulReal}
    {Df Dg : MulTangentMap}
    (hf : HasMulDerivativeAt f x Df)
    (hDf : Df.ScaleControlled)
    (hadm : ChainAdmissibleAt f x)
    (hg : HasMulDerivativeAt g (f x) Dg)
    (hDg : Dg.ScaleControlled) :
    HasMulDerivativeAt
      (fun y => g (f y))
      x
      (MulTangentMap.comp Dg Df) := by
\end{leancode}

\begin{proof}
Fix $h$ approaching the pivot and write $\name{inner} = \name{response}(f,x,h)$.

By Theorem~\ref{Analysis:Chain:thm:resp-conv}, $\name{inner}$ approaches the
pivot, so it is a legitimate perturbation and the hypothesis on $g$ may be
applied \emph{along it}.  That gives the outer error
$\name{ratio}\bigl(\name{response}(g,f(x),\name{inner})_n,\, D_g(\name{inner}_n)\bigr)$
negligible relative to $\name{inner}$; chain admissibility re-bases it against
$h$.

Separately, the hypothesis on $f$ gives the inner error
$\name{ratio}(\name{inner}_n,\, D_f(h_n))$ negligible relative to $h$.  Push it
through $D_g$ using $\name{negligible\_map}$ of
\texttt{PrimeTensor/Analysis/Control.lean} --- this is where control of $D_g$
is used --- and rewrite the image with
Lemma~\ref{Analysis:Chain:lem:map-ratio}, obtaining
$\name{ratio}\bigl(D_g(\name{inner}_n),\, D_g(D_f(h_n))\bigr)$ negligible
relative to $h$.

Multiply the two negligible germs, which is legitimate by
\texttt{PrimeTensor/Analysis/Rules.lean}.  Their product telescopes: by the
cocycle identity of Lemma~\ref{Analysis:Chain:lem:ratio},
\[
  \name{ratio}\bigl(\name{resp}_n,\, D_g(\name{inner}_n)\bigr) \cdot
  \name{ratio}\bigl(D_g(\name{inner}_n),\, D_g(D_f(h_n))\bigr)
  =
  \name{ratio}\bigl(\name{resp}_n,\, D_g(D_f(h_n))\bigr).
\]
Finally Lemma~\ref{Analysis:Chain:lem:resp-comp} identifies $\name{resp}_n$
with the response of the composite, and
$\name{comp}(D_g,D_f)(h_n) = D_g(D_f(h_n))$ by $\ctor{rfl}$, so the product is
exactly the error of the composite against its model.
\end{proof}

\begin{designnote}[A telescope where the classical proof has a sum]
\label{Analysis:Chain:note:telescope}
The classical chain rule splits the error into two pieces,
\begin{align*}
  g(f(x+h)) - g(f(x)) - D_g D_f h
  &= \bigl[g(f(x+h)) - g(f(x)) - D_g \Delta\bigr]
   + D_g\bigl[\Delta - D_f h\bigr], \\
  \Delta &= f(x+h) - f(x),
\end{align*}
and estimates the two summands separately, the second using linearity and
boundedness of $D_g$.  The proof above has the same two pieces and combines
them differently: they are ratios, they are \emph{multiplied}, and the
combination is the cocycle identity, which is exact.  There is no triangle
inequality and no constant --- the intermediate term $D_g(\name{inner}_n)$
cancels identically rather than being bounded away.

Each of the three hypotheses is spent in exactly one place, and the proof makes
that visible: control of $D_f$ only to know the inner response is a genuine
perturbation; control of $D_g$ only to push the inner error through it;
admissibility only to re-base the outer error.  The docstring says as much.
\end{designnote}

\begin{remark}[The rule cannot yet be iterated]
\label{Analysis:Chain:rem:no-iterate}
Applying Theorem~\ref{Analysis:Chain:thm:chain} twice requires knowing that
$\name{comp}(D_g,D_f)$ is itself scale controlled, and that is proved nowhere:
\texttt{PrimeTensor/Analysis/Control.lean} closed the property under product
and inverse, but composition was not defined until this file, and no lemma here
adds it.  Likewise $\name{ChainAdmissibleAt}$ is not shown closed under
composition.  So the rule is stated for one composite and has no corollary for
a tower of them.
\end{remark}

\begin{remark}[Only one function is known to be admissible]
\label{Analysis:Chain:rem:only-identity}
Lemma~\ref{Analysis:Chain:lem:adm-id} is the sole verification of
$\name{ChainAdmissibleAt}$ in the file.  Nothing establishes it for a constant
map, for a power, or for a product --- so the only $f$ with which
Theorem~\ref{Analysis:Chain:thm:chain} can currently be discharged is the
identity, where it says nothing.  This continues the pattern noted since
\texttt{PrimeTensor/Analysis/Derivative.lean}: the machinery is complete and
its nontrivial instances are still ahead.
\end{remark}

\begin{remark}[What is not proved]\label{Analysis:Chain:rem:missing}
Admissibility is not shown to follow from differentiability and control, nor
shown independent of them.  There is no uniqueness of derivatives, so
``the'' derivative of a composite is still not a phrase the development can
use.  $\name{comp}$ is given no laws --- not associativity, and not that
$\name{identity}$ is a unit for it --- and no relation is stated between
$\name{comp}$ and the pointwise product of
\texttt{PrimeTensor/Analysis/Rules.lean}.
\end{remark}

\subsection*{Summary of the correspondence}

\begin{center}
\small
\begin{tabular}{@{}lll@{}}
\hline
Lean declaration & Kind & Here \\
\hline
\lean{MulReal.ratio\_comp}          & theorem    & Lem.~\ref{Analysis:Chain:lem:ratio} \\
\lean{MulReal.ratio\_mul\_right}    & theorem    & Lem.~\ref{Analysis:Chain:lem:ratio} \\
\lean{MulReal.mul\_ratio\_right}    & theorem    & Lem.~\ref{Analysis:Chain:lem:ratio} \\
\lean{MulReal.converges\_congr}     & theorem    & Lem.~\ref{Analysis:Chain:lem:congr} \\
\lean{MulTangentMap.comp}           & definition & Def.~\ref{Analysis:Chain:def:comp} \\
\lean{MulTangentMap.comp\_apply}    & simp thm   & Def.~\ref{Analysis:Chain:def:comp} \\
\lean{MulTangentMap.map\_ratio}     & theorem    & Lem.~\ref{Analysis:Chain:lem:map-ratio} \\
\lean{MulDifferential.negligible\_congr} & theorem & Lem.~\ref{Analysis:Chain:lem:congr} \\
\lean{negligible\_converges}        & theorem    & Lem.~\ref{Analysis:Chain:lem:neg-conv} \\
\lean{response\_converges}          & theorem    & Thm.~\ref{Analysis:Chain:thm:resp-conv} \\
\lean{ChainAdmissibleAt}            & definition & Def.~\ref{Analysis:Chain:def:admissible} \\
\lean{chainAdmissibleAt\_identity}  & theorem    & Lem.~\ref{Analysis:Chain:lem:adm-id} \\
\lean{response\_comp}               & theorem    & Lem.~\ref{Analysis:Chain:lem:resp-comp} \\
\lean{hasMulDerivativeAt\_comp}     & theorem    & Thm.~\ref{Analysis:Chain:thm:chain} \\
\hline
\end{tabular}
\end{center}

\noindent
Every declaration of \texttt{PrimeTensor/Analysis/Chain.lean} appears above.
The file contains no \lean{sorry}, no axiom, and no unproved named proposition.
Design notes~\ref{Analysis:Chain:note:why-hypothesis}
and~\ref{Analysis:Chain:note:telescope} and
Remarks~\ref{Analysis:Chain:rem:diff-cont},
\ref{Analysis:Chain:rem:no-iterate}, \ref{Analysis:Chain:rem:only-identity}
and~\ref{Analysis:Chain:rem:missing} are stated for the reader and are not
declarations of the file; in particular the comparison with the classical
error split is a gloss and is not formalised.
